\section{Threats to Validity}
\label{sec:threats}

\paragraph*{Construct validity.}
Our central construct is empirical training measurement cost,
$\TrainCost$. We approximate it by the sum of measured per-configuration
times in pre-existing benchmarks. This is a deliberate simplification:
real measurement cost in production includes warm-up, repetition for noise
control, configuration-derivation cost from the feature model, and
operating-system overhead, not all of which are captured by the artifacts
we use. We mitigate this by reusing benchmark data prepared by prior work
(Kaltenecker \etal~\cite{kaltenecker2019}, Pereira
\etal~\cite{pereira2020}) where hardware, workload, and noise control are
documented. We also set model-building cost aside as second-order for the
learners we use; if it had to be added, it would be additive across
strategies and would not change the cost-difference comparisons.

\paragraph*{Aggregated labels vs.\ true acquisition effort.}
Many published datasets expose a single aggregated label per configuration
(\eg the median over $5$ runs in the \dsWLI ICSE 2023 package, or
the average over $3$ runs in the performance--energy corpus), without raw
repetition-level traces. When that is the case, we approximate true
acquisition effort as $k\times\text{median or mean}$ where $k$ is the
documented number of repetitions, and we report results so labelled. This
is an approximation: it ignores warm-up overhead, retries, and any setup
phase not reflected in the per-run aggregate. We are explicit in
\Cref{sec:results} about which datasets contributed raw traces and which
contributed aggregates with $k\times$ correction.

\paragraph*{Target and cost are independent choices.}
The cost dimension and the prediction target are independent
(\Cref{subsec:empirical-cost}). In this paper we instantiate
$\mathit{cost}(c)$ as wall-clock execution time throughout, because it is
the resource most commonly reported as constrained in this literature.
That choice is operational, not fundamental: a different operator could
legitimately use cumulative output size as $\mathit{cost}(c)$ when storage
is the binding constraint (\eg an x264 transcoding service with a disk
budget), or energy on power-bound hardware. The mathematical content of
$\TrainCost$, the calibration ratio, and the Pareto comparison is
unchanged under such substitutions; only the operational interpretation
changes. We make two specific choices for this study. First, even when
the predicted property is non-time-valued (\eg x264 output size), we use
the elapsed execution time of the run that produces the label as
$\mathit{cost}(c)$: obtaining the size still requires running the
encoder, and our operational constraint is time. Second, for throughput
datasets such as H2, throughput remains a valid prediction target and can
also serve as a cost axis after orientation, for example by using inverse
throughput so that lower values correspond to cheaper acquisition. We
therefore report these cases as throughput-based cost analyses rather than
silently mixing them with elapsed-time costs. Replicating the study under
additional non-time cost dimensions is plausible and is discussed in
\Cref{sec:discussion}.

\paragraph*{Setup overhead in full-stack systems.}
For systems where per-configuration acquisition cost is dominated by
orchestration overhead (build, warm-up, idle monitoring, deployment) rather
than by the measurement itself, the cost construct becomes the system's
deployment workflow rather than the sampling strategy's choice of
configurations. We exclude such datasets from this study (see
\Cref{subsec:datasets}). Generalizing $\TrainCost$ to those settings is
plausible but is a separate paper.

\paragraph*{Choice of cost-property mapping.}
For non-time properties (\eg x264 output size), the empirical cost is taken
from elapsed encoding time rather than from the predicted property
itself. This is the realistic choice---measuring size still requires running
the encoder---but it makes the cost term and the prediction target different
quantities, which complicates interpretation. We address this by reporting
the asymmetric case (size as target, time as cost) explicitly in RQ5 and by
comparing it to the symmetric case (time as target, time as cost).

\paragraph*{Sampling bias and noise.}
Our results can be subject to randomness in stochastic strategies and to
measurement noise in the underlying benchmarks. We mitigate the first by
running each stochastic strategy $100$ times per
$(\text{system},\text{input},\text{property},n)$ tuple and reporting medians
and IQRs; we mitigate the second by reusing benchmarks whose authors
documented stability via repeated measurements (\Cref{tab:input_videos}).

\paragraph*{Choice of learner.}
We pick stepwise multiple linear regression as the primary learner because it
is the choice in Kaltenecker \etal~\cite{kaltenecker2019} and Pereira
\etal~\cite{pereira2020}, and we add CART and random forest as robustness
checks. Different learners may yield different MRE, but the cost-side of our
analysis (RQ1, RQ2) is independent of the learner. The learner only enters
RQ3--RQ5 via $\MRE$. We acknowledge that different learner choices may shift
specific Pareto frontier pairs.

\paragraph*{Internal validity of feature exclusion.}
For each predicted property, we exclude every measured-performance column
from the feature set so that the model cannot trivially learn the target. We
list the exclusions explicitly in \Cref{subsec:protocol} (\eg \texttt{time},
\texttt{size}, \texttt{usertime}, \texttt{systemtime}, \texttt{elapsedtime},
\texttt{cpu}, \texttt{frames}, \texttt{fps}, \texttt{kbs}, \texttt{etime} for
x264 size prediction). Errors in this list would invalidate the experiment,
so the exclusion list is part of the released artifact.

\paragraph*{Generalizability across systems and contexts.}
External validity is limited to the systems and inputs we use. We mitigate
this by combining ten configurable systems already used in prior work (Tier
2) with a deeper x264 design over $17$ inputs and $2$ properties (Tier 1),
and by leaving room for Tier-3 datasets when artifacts permit. We are
explicit that any RQ4/RQ5 conclusion is a conclusion about the systems and
contexts we measured, not about configurable systems in general.

\paragraph*{Generalizability across strategies.}
We use the six literature strategies of Kaltenecker
\etal~\cite{kaltenecker2019}, plus diagnostic cheap-first/expensive-first
oracles. We do not claim these exhaust the design space of cost-aware
sampling. The diagnostic oracles are intended to bound the effect that
cost-awareness could have, not to recommend a deployable strategy.

\paragraph*{Static benchmark bias.}
The benchmarks we reuse already measured all configurations, so we work from
a complete table rather than running new measurements. This is a strength for
reproducibility but a weakness for cost realism: the cost values reflect the
hardware and workload of the original experiments. Cost heterogeneity
patterns may differ on different hardware. We address this in
\Cref{sec:discussion} by recommending a pilot procedure for practitioners to
check the cost coefficient of variation on their own machine.

\paragraph*{Statistical multiple comparisons.}
We perform many statistical tests across systems, inputs, properties, and
strategies. We will use Holm--Bonferroni correction within each RQ family of
tests and report effect sizes (Vargha--Delaney $\hat A_{12}$,
Kendall's $\tau$) alongside $p$-values to keep the conclusions tied to
practical significance and not only to statistical significance.

\paragraph*{Pre-registered thresholds for hypothesis verdicts.}
The verdicts in \Cref{tab:hypothesis-verdicts} combine a statistical
test with an effect-size or practical-magnitude threshold (e.g., $5\%$
absolute deviation from $1$ for the median calibration ratio in H2b,
$\eta^2\geq 0.01$ for H1b/H2c, $\mathrm{CV}\geq 0.20$ \emph{or}
$p_{90}/p_{10}\geq 1.5$ for H1a). Different but defensible thresholds
would shift some borderline verdicts (notably H2b on $7$Z, Dune, LLVM,
Polly, where the $95\%$ bootstrap CI excludes $1$ but the median
deviation is $<\!\!5\%$, and H5b at $4.2$ percentage points). To
keep the choice auditable, the per-scope statistic, $p$-value, effect
size, and verdict are all written to \texttt{hypothesis\_tests.csv} in
each artifact folder, so a reviewer can recompute verdicts under
alternative thresholds without rerunning any experiments.
